\section{Experimental Evaluation}
\label{sec:eval}
The evaluation examines two bounded simulator paths. R13 is an exploratory
application-level protocol and data-path run; it supports only the recorded
protocol, readback, lineage, and model-gate outcomes because its exact runtime
source was not pinned. R20 is the primary current-version result: six
sequential tasks through one registered harness across two configured scenes,
with a run-specific source reconstruction kit. It supports the tested
task-to-device integration path, not configuration-reuse benefit or process
attainment. All process experiments use repository simulators. Older
scripted and regression records remain in the companion evidence package and
are not pooled with these results.

\subsection{Evidence Units and Experimental Procedure}
\label{sec:testbed}
\label{sec:procedure}
R13 uses two configured protocol/data paths in cast-film and annealing
simulators. R20 uses the registered \texttt{omp} harness in those scenes, with
three sequential task executions per scene on one simulator instance per
scene. The companion reproducibility package records the older protocol,
controller, deployment, direct-node, and scripted-recipe cases separately;
they are not part of the current primary result.

\subsection{Exploratory Two-Scene Protocol and Data Path (R13)}
\label{sec:crossscene}
The R13 report identifies its application run as v0.8.0; the exact runtime
source was not recorded. It used cast-film and annealing simulator presets.
Five protocol families---OPC UA, MQTT, Modbus TCP, Modbus
RTU over TCP, and HTTP/REST---passed read probes; bounded writes returned
matching readbacks. Six 120-s runs per scene at 1-s sampling produced
lineage-verified datasets.

For each dataset row $i$, let $m_i$ be the last MV value in its history window
and $c_i$ the first CV value in its target horizon. The gate computes
\begin{equation}
r_{\Delta}(\ell)=\operatorname{corr}(\Delta m_i,\Delta c_{i+\ell}),
\qquad \ell\in\{1,2,3\},
\label{eq:r13corr}
\end{equation}
using all rows in the dataset for which the shifted pair is defined, and
selects the lag with greatest $|r_{\Delta}(\ell)|$. It then requires
$|r_{\Delta}|\geq0.2$, one-step test NRMSE $\leq0.10$, and at least 20\%
four-step rollout NRMSE reduction over persistence. Both
post-pilot candidates failed the gates; neither produced an accepted model,
profile, recommendation, or dispatch. The overlapping 1-s windows are
serially dependent. This run therefore supports protocol/data-path execution
and rejection behavior only, not model quality or process performance.
Signal-pair selection followed an exploratory pilot, and the run-time source
diff was not preserved; the R13 results are included as secondary diagnostic
evidence, with detailed records in the reproducibility supplement.

Table~\ref{tab:r13gates} reports the selected signal pair, the full-dataset
first-difference correlation used for eligibility screening, and two model
metrics computed on held-out test rows. Both scenes failed all three criteria;
these values document the recorded rejection path and are not confirmatory
model-performance estimates. The archived run does not identify its exact
runtime source, so the NRMSE normalization implementation cannot be verified
against that run; the companion provenance notes describe the available
candidate evaluator and this limit.
\begin{table}[!t]
\centering
\caption{R13 exploratory model-gate diagnostics. The correlation column is
the full-dataset eligibility screen; NRMSE and rollout reduction use held-out
test rows. Both candidates fail all three criteria.}
\label{tab:r13gates}
\footnotesize
\setlength{\tabcolsep}{2.2pt}
\begin{tabularx}{\columnwidth}{@{}lXccc@{}}
\toprule
Scene & MV $\to$ CV & $|r_{\Delta}|$ (lag) & 1-step NRMSE & 4-step reduction \\
\midrule
Film & \shortstack[l]{\texttt{ScrewSpeedSP}\\$\to$ \texttt{thick}} & 0.0949 (2) & 0.2587 & 14.11\% \\
Anneal & \shortstack[l]{Z3 setpoint\\$\to$ \texttt{StripTemp}} & 0.0526 (3) & 0.3632 & 4.83\% \\
\bottomrule
\end{tabularx}
\end{table}
\subsection{Registered-Harness Process Path}
\label{sec:harness-e2e}
R20 ran the isolated v0.8.0 working-tree application with the registered
\texttt{omp} adapter and GLM-5.3-Flash in cast-film and annealing simulator
scenes. The app and simulator diffs, base revisions, dependency lockfiles, and
runtime prompt assets are included with the run record; package manifests are
recovered from the recorded base revisions. The external model service does
not provide a deterministic replay of its selection. Each execution created
a separate line, product, recipe, run, channel, agent, and task context; three
sequential executions used one simulator
instance per scene, so they are not independent replicates. Cast film used
OPC UA for screw-speed control and MQTT for the thickness probe; annealing
used Modbus TCP for both the zone-3 setpoint and strip-temperature probe.
The cast-film candidate set was $\{125,130,135\}$~rpm and the annealing
candidate set was $\{735,740,745\}\,^{\circ}$C. Given a live control read,
process measurement, and target band, GLM-5.3-Flash proposed one bounded
\texttt{recipe\_trial}. A synthetic administrator approved each proposal;
the route log recorded one trial and a 1/1 verified device receipt, and all
six tasks reached the completed state. Following task completion and its
evidence queries, a 20-s wait was initiated; one driver probe per task was
then ingested through the isolated SQLite time-series fallback and returned
under the active run tag. The integration checks cover per-execution context, an approved bounded
proposal, route-specific receipt and readback, task completion, and an
active-run-tagged sample; all six records contain this software path. The
runner recorded 76 passing checks and no failures. The task history contains
mirrored event projections, so the per-task operation count is taken from the
route-specific recipe and ACK records.

Table~\ref{tab:q8} reports the pre-task measurement, selected setpoint, and
single post-trial sample. All six post-trial values remained outside the
declared target bands. This experiment verifies the bounded
harness-to-simulator execution and evidence path; it does not show process
attainment, convergence, or quality improvement. The single sample is not a
response trajectory, and the SQLite fallback does not test the normal online
DAQ backend. The two-scene run also does not isolate profile/binding changes
or establish cross-harness portability.

\begin{table}[!t]
\centering
\caption{Model-selected settings and simulator measurements. Each row
reports the pre-task value, configured target band, manipulated variable's
initial-to-selected setting, and one run-tagged sample taken after a 20-s wait
initiated following task completion and evidence queries. All six samples are
outside target; executions are sequential within each scene.}
\label{tab:q8}
\footnotesize
\setlength{\tabcolsep}{2.5pt}
\begin{tabularx}{\columnwidth}{@{}lXcX@{}}
\toprule
Case & Pre-task CV / target & Setting: start$\to$trial & CV after wait \\
\midrule
Film-1 & $55.82/48$--$52~\mu$m & $120\to125$~rpm & $45.69~\mu$m \\
Film-2 & $45.69/48$--$52~\mu$m & $120\to130$~rpm & $46.69~\mu$m \\
Film-3 & $46.69/48$--$52~\mu$m & $120\to125$~rpm & $45.69~\mu$m \\
Anneal-1 & $675.40/700$--$760~^{\circ}$C & $730\to745~^{\circ}$C & $679.60~^{\circ}$C \\
Anneal-2 & $678.30/700$--$760~^{\circ}$C & $730\to745~^{\circ}$C & $677.70~^{\circ}$C \\
Anneal-3 & $676.90/700$--$760~^{\circ}$C & $730\to745~^{\circ}$C & $678.10~^{\circ}$C \\
\bottomrule
\end{tabularx}
\end{table}

Earlier protocol and admission-ablation suites, deterministic-controller
runs, deployment regressions, direct-node missions, and a scripted recipe
case are retained in the companion reproducibility package for auditability.
They use older versions or interfaces and are not pooled with R13 or R20, used
as matched strategy comparisons, or relied on for the current integration
claims.
